% =============================================================================
% NER EXTENSION: ENTITY-AWARE QUESTION GENERATION AND ANSWERING
% Template: IEEEtran.cls 2015/08/26 version V1.8b
% =============================================================================
% This file contains the Methodology and Experiments sections for the 
% NER-extension component of the AskQE paper applied to BioMQM.
% Include this in your main .tex file after the appropriate \section{} commands
% =============================================================================

% -----------------------------------------------------------------------------
% METHODOLOGY SECTION
% -----------------------------------------------------------------------------

\subsection{Entity-Aware Question Generation}
\label{sec:ner_methodology}

The baseline AskQE pipeline generates questions from a source sentence $s$ using a vanilla prompting strategy, where the language model produces a set of questions $\mathcal{Q} = \{q_1, \ldots, q_k\}$ conditioned solely on the surface form of the sentence. While effective, this approach may fail to target domain-specific entities that are critical for assessing translation quality in specialised domains such as biomedical text. We propose an \textit{entity-aware} extension that leverages Named Entity Recognition (NER) to guide question generation toward salient biomedical concepts.

\subsubsection{NER-Guided Pipeline Overview}

The entity-aware extension augments the standard AskQE pipeline with a three-stage process: (1) biomedical named entity extraction, (2) entity-conditioned question generation, and (3) entity-specific question answering. The key intuition is that translation errors affecting domain-critical entities---such as drug names, diseases, or anatomical structures---are more likely to indicate clinically significant quality degradation.

\subsubsection{Biomedical Named Entity Extraction}

Given a source sentence $s$, we extract a set of biomedical entities $\mathcal{E} = \{e_1, \ldots, e_m\}$ using a dedicated NER model. Each entity $e_j$ is represented as a tuple:
\begin{equation}
    e_j = (\text{text}_j, \, \tau_j, \, p_j, \, \text{start}_j, \, \text{end}_j)
\end{equation}
where $\text{text}_j$ is the entity surface form, $\tau_j \in \mathcal{T}$ denotes its semantic type, $p_j \in [0, 1]$ is the prediction confidence, and $(\text{start}_j, \text{end}_j)$ are the character-level span offsets.

We employ the \texttt{d4data/biomedical-ner-all} model~\cite{d4data2022}, a transformer-based token classifier fine-tuned on multiple biomedical NER corpora. The model supports a rich type inventory $\mathcal{T}$ that includes, among others:
\begin{itemize}
    \item \textsc{Disease}: pathological conditions and disorders
    \item \textsc{Drug}: pharmaceutical substances and medications
    \item \textsc{Symptom}: clinical signs and symptoms
    \item \textsc{Gene}: gene and protein names
    \item \textsc{Detailed\_description}: anatomical and procedural descriptions
\end{itemize}

To handle subword tokenisation artefacts, we apply an aggregation strategy that merges consecutive entities of the same type. For two adjacent entities $e_j$ and $e_{j+1}$ satisfying:
\begin{equation}
    \text{start}_{j+1} \leq \text{end}_j + 1 \quad \land \quad \tau_{j+1} = \tau_j
\end{equation}
we merge them into a single entity with concatenated text and averaged confidence:
\begin{equation}
    p_{\text{merged}} = \frac{p_j + p_{j+1}}{2}
\end{equation}

Additionally, entities with surface forms shorter than two characters are discarded as likely tokenisation noise, and entities tagged with the non-entity label (\texttt{O}) are filtered out.

\subsubsection{Entity-Conditioned Question Generation}

For each extracted entity $e_j$, we generate a question $q_j$ that specifically targets the information conveyed by that entity in the context of the source sentence. Unlike the baseline approach, which generates questions about the sentence as a whole, entity-conditioned generation constrains the question's focus to a specific biomedical concept.

The question generation is formulated as a conditional language generation task. Given the sentence $s$, the entity text $\text{text}_j$, and the entity type $\tau_j$, we prompt Qwen2.5-3B-Instruct with a structured template:
\begin{equation}
    q_j = \text{LLM}\big(s, \, \text{text}_j, \, \tau_j \,;\, \theta_{\text{QG}}\big)
\end{equation}

The prompt instructs the model to generate exactly one clear, specific question such that: (i) the question can be answered using only information present in the sentence, (ii) the question specifically asks about the designated entity, and (iii) the expected answer is the entity itself or information directly related to it. The generation uses a low temperature ($T = 0.3$) to encourage deterministic, focused questions.

This yields a set of entity-specific questions $\mathcal{Q}_{\text{NER}} = \{q_1, \ldots, q_m\}$, where $|\mathcal{Q}_{\text{NER}}| = |\mathcal{E}|$. A fundamental difference from the baseline is that $|\mathcal{Q}_{\text{NER}}|$ varies per sentence depending on the number of extracted entities, whereas the baseline generates a fixed number of general-purpose questions.

\subsubsection{Entity-Specific Question Answering}

Each entity-specific question $q_j$ is answered in two contexts: the source sentence $s$ and the backtranslated machine translation output $s_{\text{bt}}$:
\begin{align}
    a_j^{(s)} &= \text{LLM}(s, \, q_j \,;\, \theta_{\text{QA}}) \\
    a_j^{(bt)} &= \text{LLM}(s_{\text{bt}}, \, q_j \,;\, \theta_{\text{QA}})
\end{align}

The answering model (Qwen2.5-3B-Instruct) is instructed to respond based on the sentence alone, returning \texttt{[NOT FOUND]} if the answer is not present. This extractive constraint ensures that answer discrepancies between $a_j^{(s)}$ and $a_j^{(bt)}$ reflect genuine translation divergences rather than generation artefacts.

\subsubsection{Answer Comparison and Scoring}

For each question $q_j$, translation quality is assessed by comparing the source and backtranslated answers using multiple evaluation metrics. Let $\phi : \mathcal{A} \times \mathcal{A} \rightarrow \mathbb{R}$ denote a generic comparison function. We compute:

\paragraph{Token-level F1.}
\begin{equation}
    \text{F1}(a^{(s)}_j, a^{(bt)}_j) = \frac{2 \cdot |\,\text{tokens}(a^{(s)}_j) \cap \text{tokens}(a^{(bt)}_j)\,|}{|\text{tokens}(a^{(s)}_j)| + |\text{tokens}(a^{(bt)}_j)|}
\end{equation}

\paragraph{Exact Match.}
\begin{equation}
    \text{EM}(a^{(s)}_j, a^{(bt)}_j) = \mathbb{1}\big[\text{norm}(a^{(s)}_j) = \text{norm}(a^{(bt)}_j)\big]
\end{equation}
where $\text{norm}(\cdot)$ applies whitespace normalisation and lowercasing.

\paragraph{BLEU and chrF.} We additionally compute sentence-level BLEU~\cite{papineni2002bleu} and character-level F-score (chrF)~\cite{popovic2015chrf} to capture $n$-gram and character-level overlap, respectively.

\paragraph{SBERT Cosine Similarity.} We employ Sentence-BERT~\cite{reimers2019sentence} to compute the embedding-based semantic similarity:
\begin{equation}
    \text{SBERT}(a^{(s)}_j, a^{(bt)}_j) = \frac{\mathbf{h}_{a^{(s)}_j} \cdot \mathbf{h}_{a^{(bt)}_j}}{\|\mathbf{h}_{a^{(s)}_j}\| \cdot \|\mathbf{h}_{a^{(bt)}_j}\|}
\end{equation}
where $\mathbf{h}_a \in \mathbb{R}^d$ is the sentence embedding of answer $a$.

\paragraph{Sentence-Level Aggregation.}
For a sentence $s_i$ with $m_i$ extracted entities, the overall sentence-level score under metric $\phi$ is computed as the average across all entity-specific comparisons:
\begin{equation}
    \Phi(s_i) = \frac{1}{m_i} \sum_{j=1}^{m_i} \phi\big(a_j^{(s)}, \, a_j^{(bt)}\big)
    \label{eq:sentence_agg}
\end{equation}


% -----------------------------------------------------------------------------
% EXPERIMENTS & RESULTS SECTION
% -----------------------------------------------------------------------------

\subsection{NER Extension Results on BioMQM}
\label{sec:ner_results}

We evaluate the entity-aware extension on the BioMQM dataset, comparing it against the baseline AskQE pipeline and three prompt ablation variants: (P1) few-shot, (P2) chain-of-thought, and (P3) concise prompting. All configurations employ Qwen2.5-3B-Instruct as the underlying language model, ensuring a controlled comparison where the only variable is the question generation strategy.

\subsubsection{Overall Performance by Language}

Table~\ref{tab:ner_by_language} presents the average evaluation metrics aggregated by language across all severity levels.

\begin{table*}[htbp]
\centering
\caption{Evaluation metrics by language: baseline vs.\ NER extension vs.\ prompt ablation variants on BioMQM. Deltas ($\Delta$) are computed with respect to the baseline. Best scores per language are \textbf{bolded}.}
\label{tab:ner_by_language}
\resizebox{\textwidth}{!}{%
\begin{tabular}{l|ccccc|ccccc}
\hline
& \multicolumn{5}{c|}{\textbf{F1}} & \multicolumn{5}{c}{\textbf{Exact Match}} \\
\textbf{Lang} & Baseline & NER & P1 & P2 & P3 & Baseline & NER & P1 & P2 & P3 \\
\hline
de & 0.517 & \textbf{0.675} & 0.591 & 0.513 & 0.522 & 0.360 & \textbf{0.549} & 0.291 & 0.339 & 0.305 \\
es & 0.453 & \textbf{0.686} & 0.563 & 0.472 & 0.491 & 0.243 & \textbf{0.512} & 0.223 & 0.244 & 0.259 \\
fr & 0.434 & \textbf{0.640} & 0.529 & 0.467 & 0.481 & 0.265 & \textbf{0.433} & 0.230 & 0.280 & 0.261 \\
ru & 0.484 & \textbf{0.710} & 0.548 & 0.490 & 0.516 & 0.255 & \textbf{0.481} & 0.164 & 0.230 & 0.254 \\
zh-CN & 0.419 & \textbf{0.608} & 0.518 & 0.449 & 0.458 & 0.196 & \textbf{0.404} & 0.149 & 0.199 & 0.193 \\
\hline
\end{tabular}%
}

\vspace{0.3cm}

\resizebox{\textwidth}{!}{%
\begin{tabular}{l|ccccc|ccccc}
\hline
& \multicolumn{5}{c|}{\textbf{chrF}} & \multicolumn{5}{c}{\textbf{BLEU}} \\
\textbf{Lang} & Baseline & NER & P1 & P2 & P3 & Baseline & NER & P1 & P2 & P3 \\
\hline
de & 59.17 & 59.17 & \textbf{60.87} & 53.22 & 54.87 & 45.71 & \textbf{47.97} & 43.25 & 38.36 & 39.49 \\
es & 51.43 & \textbf{55.84} & 57.45 & 50.33 & 51.81 & 36.97 & \textbf{45.24} & 39.41 & 34.17 & 35.73 \\
fr & 50.76 & \textbf{58.01} & 55.21 & 48.75 & 51.33 & 36.36 & \textbf{43.48} & 37.43 & 33.38 & 36.98 \\
ru & 54.52 & \textbf{62.32} & 57.10 & 51.69 & 53.97 & 38.45 & \textbf{48.47} & 34.97 & 33.38 & 37.32 \\
zh-CN & 48.57 & \textbf{55.69} & 54.55 & 47.09 & 48.73 & 32.50 & \textbf{42.37} & 32.71 & 29.39 & 31.21 \\
\hline
\end{tabular}%
}
\end{table*}

The NER extension yields consistent and substantial improvements over the baseline across all five languages and all evaluation metrics. The most pronounced gains are observed for F1 and Exact Match: the average F1 improvement ranges from $+0.158$ (de) to $+0.233$ (es), while Exact Match improves by $+0.168$ (fr) to $+0.269$ (es). These results represent relative improvements of approximately 30--50\% over the baseline, demonstrating the effectiveness of entity-conditioned question generation.

Notably, the NER extension outperforms all three prompt ablation variants (P1, P2, P3) by a wide margin. While the few-shot variant (P1) provides modest improvements over the baseline in F1 (averaging $+0.09$), the NER extension surpasses it by an additional $+0.11$ on average. The chain-of-thought (P2) and concise (P3) variants offer negligible or marginal improvements over the baseline, and in some cases even decrease performance (e.g., P2 chrF for German: $-5.95$), suggesting that prompt engineering alone does not address the fundamental limitation of question specificity.

\subsubsection{Semantic Similarity Analysis (SBERT)}

Table~\ref{tab:ner_sbert} reports the SBERT cosine similarity scores by language, providing an embedding-based perspective on answer agreement.

\begin{table}[htbp]
\centering
\caption{Average SBERT cosine similarity by language. $\Delta$ denotes the improvement over the baseline.}
\label{tab:ner_sbert}
\begin{tabular}{l|cc|cccc}
\hline
\textbf{Lang} & \textbf{Baseline} & \textbf{NER} & \textbf{P1} & \textbf{P2} & \textbf{P3} & $\boldsymbol{\Delta}_{\textbf{NER}}$ \\
\hline
de & 0.661 & \textbf{0.766} & 0.724 & 0.642 & 0.656 & $+$0.105 \\
es & 0.611 & \textbf{0.750} & 0.702 & 0.608 & 0.625 & $+$0.139 \\
fr & 0.607 & \textbf{0.711} & 0.691 & 0.620 & 0.636 & $+$0.103 \\
ru & 0.631 & \textbf{0.798} & 0.713 & 0.633 & 0.647 & $+$0.168 \\
zh-CN & 0.588 & \textbf{0.714} & 0.702 & 0.611 & 0.611 & $+$0.125 \\
\hline
\end{tabular}
\end{table}

The SBERT results corroborate the findings from lexical metrics, with the NER extension achieving the highest semantic similarity across all languages. The improvements are particularly striking for Russian ($+0.168$), which is the most morphologically complex language in our evaluation, and for Spanish ($+0.139$). The consistency of improvements across both lexical and embedding-based metrics suggests that entity-awareness enhances not merely surface-level token overlap, but also the deeper semantic alignment between source and backtranslated answers.

Interestingly, the P2 (chain-of-thought) and P3 (concise) variants yield SBERT scores comparable to or even marginally below the baseline (e.g., P2-de: $0.642$ vs.\ baseline: $0.661$), indicating that more elaborate prompting does not inherently improve semantic consistency when questions remain untargeted.

\subsubsection{Analysis by Error Severity}

Table~\ref{tab:ner_severity} presents the F1 and SBERT scores broken down by MQM error severity, allowing us to assess whether the NER extension provides differential benefits depending on the severity of the translation error.

\begin{table*}[htbp]
\centering
\caption{F1 and SBERT scores by language and error severity. $\Delta$ denotes the improvement of NER over baseline.}
\label{tab:ner_severity}
\resizebox{\textwidth}{!}{%
\begin{tabular}{ll|cc|c|cc|c}
\hline
& & \multicolumn{3}{c|}{\textbf{F1}} & \multicolumn{3}{c}{\textbf{SBERT}} \\
\textbf{Lang} & \textbf{Severity} & Baseline & NER & $\boldsymbol{\Delta}$ & Baseline & NER & $\boldsymbol{\Delta}$ \\
\hline
\multirow{4}{*}{de}
& Neutral & 0.604 & 0.705 & $+$0.101 & 0.749 & 0.767 & $+$0.017 \\
& Minor & 0.526 & 0.719 & $+$0.193 & 0.663 & 0.802 & $+$0.139 \\
& Major & 0.490 & 0.649 & $+$0.160 & 0.637 & 0.752 & $+$0.115 \\
& Critical & 0.533 & 0.796 & $+$0.263 & 0.648 & 0.863 & $+$0.216 \\
\hline
\multirow{4}{*}{es}
& Neutral & 0.527 & 0.673 & $+$0.146 & 0.667 & 0.722 & $+$0.055 \\
& Minor & 0.441 & 0.698 & $+$0.258 & 0.599 & 0.763 & $+$0.163 \\
& Major & 0.444 & 0.597 & $+$0.153 & 0.590 & 0.699 & $+$0.109 \\
& Critical & 0.430 & 0.615 & $+$0.184 & 0.634 & 0.682 & $+$0.047 \\
\hline
\multirow{4}{*}{fr}
& Neutral & 0.522 & 0.640 & $+$0.118 & 0.692 & 0.702 & $+$0.010 \\
& Minor & 0.393 & 0.647 & $+$0.254 & 0.565 & 0.719 & $+$0.154 \\
& Major & 0.391 & 0.639 & $+$0.248 & 0.574 & 0.724 & $+$0.150 \\
& Critical & 0.198 & 0.206 & $+$0.009 & 0.453 & 0.391 & $-$0.062 \\
\hline
\multirow{4}{*}{ru}
& Neutral & 0.498 & 0.709 & $+$0.212 & 0.653 & 0.800 & $+$0.148 \\
& Minor & 0.500 & 0.744 & $+$0.244 & 0.624 & 0.822 & $+$0.198 \\
& Major & 0.375 & 0.678 & $+$0.303 & 0.511 & 0.787 & $+$0.275 \\
& Critical & 0.333 & 0.481 & $+$0.148 & 0.499 & 0.541 & $+$0.042 \\
\hline
\multirow{4}{*}{zh-CN}
& Neutral & 0.433 & 0.648 & $+$0.215 & 0.605 & 0.766 & $+$0.161 \\
& Minor & 0.422 & 0.588 & $+$0.166 & 0.583 & 0.683 & $+$0.100 \\
& Major & 0.379 & 0.600 & $+$0.221 & 0.574 & 0.718 & $+$0.144 \\
& Critical & 0.453 & 0.586 & $+$0.133 & 0.663 & 0.677 & $+$0.015 \\
\hline
\end{tabular}%
}
\end{table*}

Several key observations emerge from the severity-stratified analysis:

\paragraph{Broad Improvements Across All Severities.}
The NER extension improves performance across all severity levels in virtually every language--metric combination. The improvements are not confined to a single severity category, indicating that entity-aware question generation provides a fundamentally superior signal for translation quality assessment regardless of the nature of the underlying error.

\paragraph{Largest Gains on Minor and Major Errors.}
The most substantial F1 improvements typically occur for Minor and Major severity errors. For example, in Russian, Major-severity F1 improves by $+0.303$ (from~$0.375$ to $0.678$), and the corresponding SBERT improvement is $+0.275$. Similarly, French Minor-severity F1 improves by $+0.254$. These results suggest that entity-focused questions are particularly effective at detecting errors that affect specific biomedical terms---the types of errors most likely to fall within Minor and Major severity categories.

\paragraph{Strong Performance on Critical Errors.}
For Critical errors, the NER extension generally maintains improvements, albeit with more variability due to smaller sample sizes. German Critical-severity F1 improves by $+0.263$ (from $0.533$ to $0.796$), while Russian Critical-severity shows a more modest $+0.148$. The one exception is French Critical severity ($n = 11$), where both approaches yield low scores (baseline $0.198$, NER $0.206$), with the SBERT metric slightly declining ($-0.062$). This anomalous result is likely attributable to the very small sample size and the potential presence of complex, multi-error sentences in this stratum.

\paragraph{Neutral Severity as a Baseline Check.}
The Neutral severity category---representing sentences without translation errors---serves as a sanity check. Here, the NER extension consistently achieves higher scores than the baseline (e.g., de: $0.604 \rightarrow 0.705$), confirming that entity-aware questions produce more consistent answers when source and backtranslated sentences are semantically equivalent. This can be interpreted as a reduction in false positives: the baseline's more generic questions occasionally yield spurious answer differences even for correct translations, whereas entity-targeted questions are more focused and thus less prone to such noise.

\subsubsection{Discussion}

The substantial and consistent improvements delivered by the NER extension can be attributed to several complementary factors:

\begin{enumerate}
    \item \textbf{Targeted questioning.} By anchoring each question to a specific biomedical entity, the NER extension generates questions that probe precisely the information most likely to be affected by translation errors. Generic questions may inadvertently target non-essential parts of the sentence, diluting the quality signal.
    
    \item \textbf{Domain-aware entity grounding.} The use of a biomedical NER model ensures that the extracted entities are clinically relevant (e.g., drug names, diseases, symptoms), aligning the evaluation with the types of errors that matter most in medical translation contexts.
    
    \item \textbf{Scalable question cardinality.} Unlike the baseline, which generates a fixed number of questions regardless of sentence complexity, the NER extension adapts the number of questions to the information density of the sentence. Sentences rich in biomedical entities receive more probing questions, while simpler sentences receive fewer, yielding a more efficient allocation of the evaluation budget.
    
    \item \textbf{Superiority over prompt engineering.} The comparison with prompt ablation variants (P1--P3) demonstrates that the improvements are not achievable through prompt engineering alone. Few-shot examples (P1), chain-of-thought reasoning (P2), and concise instructions (P3) modify \textit{how} questions are generated but not \textit{what} they are generated about. The NER extension addresses this latter dimension by providing explicit entity-level guidance.
\end{enumerate}

\subsubsection{Limitations}

Despite the strong empirical results, several limitations should be noted:

\begin{itemize}
    \item The NER model (\texttt{d4data/biomedical-ner-all}) was trained predominantly on English biomedical text. While it performs well on English source sentences, its entity extraction quality may degrade on non-English text or on source sentences with unusual domain terminology.
    
    \item The number of generated questions is directly tied to the number of extracted entities. For sentences with few or no recognised entities, the extension may produce fewer questions than the baseline, potentially reducing evaluation coverage for such cases.
    
    \item The entity aggregation heuristic (merging consecutive same-type entities) is a rule-based post-processing step and may occasionally produce incorrect merges, particularly for complex biomedical nomenclature with nested entity boundaries.
    
    \item The evaluation on Critical severity is limited by small sample sizes in several language pairs (e.g., $n = 11$ for French, $n = 13$ for German), making the corresponding estimates less reliable and more sensitive to individual outliers.
\end{itemize}
